\documentclass{article}
\usepackage{amsmath}
\begin{document}

\section{Keys}

A plain paragraph of English words that is long enough to wrap onto several lines in the editor pane, with a comma, a semicolon; a colon: and (some parentheses) so that word and line motion have something to cross. The caret walks over it again and again.

Inline math $a^2+b^2=c^2$ sits between words, and a citation \cite{knuth} or a reference \ref{sec:keys} follows the formula $x_i$ in the same paragraph.

\begin{equation}
E = mc^2
\end{equation}

The paragraph after a display formula starts here and carries on for a few more words to fill the line.

\begin{itemize}
\item The first list item has enough words to be a real item.
\item The second list item is here too.
\end{itemize}

\begin{tabular}{ll}
alpha beta & gamma delta \\
epsilon zeta & eta theta \\
\end{tabular}

中文段落包含汉字和标点符号，用来测试按词移动和删除。日本語の段落にはひらがな、カタカナ、漢字が含まれています。한국어 문단도 있습니다.

Punctuation: well-known, state-of-the-art "quoted words", it's (nested [brackets]) and a/b paths... end.

Emoji 👍🏽 and a family 👨‍👩‍👧 and a flag 🇫🇷 and combining é and Hangul 한글 and Devanagari हिन्दी text.

\begin{verbatim}
code line one
code line two
\end{verbatim}

\end{document}
